\chapter{Objetivos}
\label{sec:objetivos}

Este capítulo apresenta os objetivos que orientam o desenvolvimento da pesquisa.

\section{Objetivo Geral}

O objetivo geral deste trabalho é [verbo no infinitivo -- por exemplo: analisar, desenvolver, propor, comparar] [o que será feito], por meio de [abordagem/metodologia empregada], considerando [principais variáveis ou aspectos avaliados]\footnote{Exemplo de nota de rodapé com referência a um site: Disponível em: \url{https://www.exemplo.com/}. Acesso em: DD mês. AAAA.}.

\section{Objetivos Específicos}

OBS: Não confunda objetivo específico com metodologia de trabalho, isso é bastante comum.

Os objetivos específicos deste trabalho são:

\begin{itemize}[leftmargin=*,itemsep=0.5em]

\item Primeiro objetivo específico, geralmente relacionado à etapa inicial ou preparatória da pesquisa (ex.: levantar, revisar, implementar);

\item Segundo objetivo específico, relacionado à execução central da pesquisa (ex.: comparar, avaliar, aplicar);

\item Terceiro objetivo específico, relacionado à análise dos dados obtidos (ex.: analisar o impacto de determinada variável sobre os resultados);

\item Quarto objetivo específico, relacionado à conclusão ou síntese dos resultados (ex.: identificar a solução mais adequada segundo os critérios definidos).

\end{itemize}
